% The command path, with the gap drawn where it is.
% Labels on short segments are placed as explicit nodes: a node[lbl] on a
% segment shorter than its own text runs into whatever is beside it.
\begin{tikzpicture}[font=\sffamily\small, x=1cm, y=1cm]

% ---------------- the node's real output path
\node[fw, text width=32mm] (loop) at (-5.9,3.6)
  {the period handler\\{\scriptsize chapter 2, and the controller in chapter 16}};
\node[hw, text width=32mm] (tim) at (-5.9,1.4)
  {advanced timer\\{\scriptsize centre aligned, three complementary pairs}};
\node[hw, text width=32mm] (dt) at (-5.9,-1.2)
  {dead-time generator\\{\scriptsize piecewise encoding, measured in step 5}};
\node[hw, text width=32mm] (pins) at (-5.9,-3.6)
  {six output pins\\{\scriptsize real signals, and the only thing here an oscilloscope would see}};
\begin{scope}[on background layer]
\node[layer, fit=(loop)(tim)(dt)(pins), inner sep=6pt,
      label={[layerlabel, anchor=south west]north west:real, and measured}] {};
\end{scope}

% ---------------- the break path, and what is absent
\node[hw, text width=26mm] (brk) at (-0.7,-1.2)
  {break input\\{\scriptsize drives the outputs inactive in hardware}};
\node[absentblk, text width=28mm] (drv) at (-0.7,-3.6)
  {gate driver, bridge, motor, current sensing\\{\scriptsize none of it on this bench}};

% ---------------- the model
\node[modelled, text width=32mm] (plant) at (4.4,-1.2)
  {plant model\\{\scriptsize inertia, damping, a crude friction term. Four parameters, none measured}};
\node[modelled, text width=32mm] (torque) at (4.4,-3.6)
  {torque\\{\scriptsize modelled, because there is no current loop}};

% ---------------- and back through real hardware
\node[hw, text width=32mm] (gen) at (4.4,1.4)
  {quadrature generator\\{\scriptsize chapter 5, now driven by the plant}};
\node[hw, text width=32mm] (dec) at (4.4,3.6)
  {the real decoder\\{\scriptsize chapter 5: quantisation, extension, index}};
\begin{scope}[on background layer]
\node[layer, fit=(gen)(dec), inner sep=6pt,
      label={[layerlabel, anchor=south west]north west:real again}] {};
\end{scope}

% ---------------- edges
\draw[flow] (loop) -- node[lbl, right]{duty} (tim);
\draw[flow] (tim) -- (dt);
\draw[flow] (dt) -- (pins);
\draw[flow] (pins) -- node[lbl, above]{would drive} (drv);
\draw[flow, draw=linecol!55] (drv) -- node[lbl, above]{so:} (plant);
\draw[flow] (plant) -- (torque);
\draw[flow] (plant) -- (gen);
\draw[flow] (gen) -- (dec);
\draw[flow] (dec) -- node[lbl, above]{position} (loop);
\draw[faultedge] (brk) -- node[lbl, above]{stops} (dt);

% the loopback, routed clear of both columns on the left
\draw[bus, draw=linecol!55] (pins.west) -- ++(-1.1,0) |- (tim.west);
\node[font=\sffamily\scriptsize, fill=white, inner sep=1.5pt, align=center]
  at (-9.4,0.1) {loopback:\\two wires,\\two captures};

\node[font=\sffamily\scriptsize, fill=white, inner sep=1.5pt, anchor=west]
  at (4.6,0.1) {velocity sets the rate};
\node[font=\sffamily\scriptsize, fill=white, inner sep=1.5pt, anchor=west]
  at (4.6,2.5) {three wires, out and back};

\node[umlnote, text width=58mm, anchor=north west] at (-8.0,-5.2)
  {A plant model read directly by the loop is a number the same processor wrote
   a microsecond earlier: no quantisation, no decode latency, no extension
   arithmetic, no index. Routed through the generator and the real decoder it
   has all four, so the controller in chapter 16 meets the signal a real encoder
   would give it. The torque is fictional; the path from command to measured
   position is not.};
\end{tikzpicture}
